\section{Introduction}
\label{sec:Introduction}
Consider the following classification situation. A university assigns students to various programs based on their scores in an aptitude test using a decision rule that maps scores to programs. For example, one such rule could be to assign the Physics program to students who have scored well in the physics section of the exam, regardless of their total score. For transparency, the university must declare its decision rule before the students write the exam. Every student taking the exam has a \textit{true aptitude} corresponding it with a \textit{true program}.
However, the students may have preferences for programs that may differ from their true programs. Knowing the decision rule, they can seek external help, such as coaching, to improve their scores, incurring some costs for such assistance. 
Now suppose the university offers two related programs (say theoretical and experimental physics programs), where the students with an aptitude in one program may prefer the other program and, therefore, prepare for the entrance exam accordingly. The university aims to create a question paper for the entrance exam  in such a manner as to assign the maximum number of students to their true programs. The challenge is that the university cannot measure the true aptitude of a student and must rely on the exam scores, i.e., its \textit{demonstrated aptitude}, to assign the program. Moreover, the university does not know the \textit{true function} (which maps true aptitude to a true program) and must learn it from past data.

The university admission problem is an example of a \textit{strategic} classification problem. The strategic classification problem extends the classification problem to the regime where, during the classifier's testing phase, a strategic data sender can manipulate feature vectors to obtain a preferred label while incurring some cost. A non-strategic or \textit{naive classifier} places the decision boundary precisely at the class boundary. Since this classifier's decision rule is known to the strategic sender, it is prone to manipulation. A \textit{strategic classifier} is robust to such manipulation and, in general, differs non-trivially from the naive classifier.

Existing works on the strategic classification problem typically assume that the sender has a uniform preference for the labels (e.g., \cite{hardt2016strategic}). The standard example of this setting is the email classification problem, where the strategic sender prefers the non-spam label over the spam label, irrespective of the underlying email content. The strategic classifier for this case is the one that shifts the decision boundary into the \textit{non-spam} region of the feature space, thereby making the criterion for classification as the non-spam label more stringent.

This chapter considers a binary classification problem where the strategic sender has a possibly non-uniform preference for labels, as illustrated in the above example for admission to theoretical and experimental physics programs. This assumption poses a new challenge as preferred feature space for one class overlaps with the space not preferred by the other. In particular, the structure of the strategic classifier obtained by \cite{hardt2016strategic} no longer holds without a universally preferred label.
The problem of constructing a strategic classifier with non-uniform preferences results in two subproblems: a \textit{decision problem} and a \textit{learning problem}. In the decision problem, the university knows the \textit{true function} (map between the true aptitude and the true program) and has to calculate the strategic classifier from it. In the learning problem, the university does not know the true function and has to learn the strategic classifier from a finite sample of the true function.

Our main findings are that classification performance improves significantly when the classifier is either allowed to make randomized decisions or allowed to \textit{not} classify certain feature vectors. Moreover, in certain problem instances, particularly those involving complex hypothesis classes, learning can be made easier by incorporating the sender's strategic behaviour in the classifier's learning phase.

\subsection{Contributions}

We adopt a game theoretic viewpoint for our analysis. We model the binary classification problem with non-uniform preferences as a screening game\footnote{In a screening game, the receiver commits to a strategy first; the sender then observes this commitment and chooses a strategy in response, after which the receiver plays the committed strategy.} between a sender and a receiver (classifier), with the receiver as the leader. The sender wants to maximize the difference between its utility and its cost and the receiver wants to maximize the probability of correctly classifying the feature vectors. The strategic classifier corresponds to the receiver's Stackelberg equilibrium strategy. For most of our results, we work with a scalar feature space. We show that under certain assumptions, the problem with the vector feature space with a linear class boundary can be reduced to the problem we consider. 

We introduce two main ideas. The first is to expand the classifier's capabilities in two ways. One way is by introducing a new penalty label $\Delta$, by which means the receiver ignores a feature vector assigned to this label. From the sender's perspective, this situation is equivalent to receiving a $-\infty$ utility for the label $\Delta$. In the context of university admission problems, the student is the sender, and the university is the receiver. The exam paper comprises of a subset of questions from a question bank, with the requirement that the student has to attempt exactly one question. The university assigns a label to each question in the exam paper with the implication that a student who correctly answers that question gets assigned that label or program. Assignment of label $\Delta$ to a question amounts to \textit{dropping} that question from the exam paper. This approach is in stark contrast to the uniform preference case, where \textit{all} the questions in the question bank are included in the question paper, and only \textit{cut-off marks} for the preferred program are increased.
Ignoring feature vectors can also be considered a generalization of feature selection.

Another way to expand the classifier's action set is by allowing randomized classification. In the university admission problem, this is analogous to saying that solving a particular question is correlated (but not with certainty) with admission to a specific program. These ideas for expanding the classifier's action set are motivated from mechanism design, where penalties influence a strategic sender's behavior, and game theory, where randomization allows for richer outcomes.

The second main idea involves incorporating the underlying game-theoretic structure of the problem by modifying the hypothesis class used for learning. Learning from the original hypothesis class will be inefficient if the university's inductive bias does not consider students' strategic nature. We show that incorporating the game's strategic structure through hypothesis class modification can make an unlearnable hypothesis class learnable.

Using these two ideas, we address the decision and learning problems of strategic classification in the scalar feature space setting. For the decision problem, our first finding is that a deterministic strategic classifier using a penalty label makes fewer classification errors than any classifier that does not use a penalty label, at least strictly fewer errors in some instances of the problem. 
With probability $1 - \delta$ (where we will write this in sections), a stochastic strategic classifier makes fewer errors than any deterministic classifier, including those that use a penalty label, and at least strictly fewer errors than a classifier with a penalty label in some instances of the problem.

However, the stochastic classifier uses only randomization on the original label set; using the penalty label does not increase its accuracy. Nonetheless, the penalty label becomes relevant when it is constrained to use a deterministic strategic classifier, a common situation encountered in real-life applications. Hence, the penalty label can be thought of as partially compensating the classifier for the absence of randomization.

The classifier described above assigns the label $\Delta$ to feature vectors near the class boundary. Indeed, multiple such classifiers exist, all with the same classification error, but they differ in the following three aspects.
The first aspect is classifier complexity, measured by the number of decision boundaries. A decision boundary is a feature vector at which the classifier changes its assigned label. Thus, the number of decision boundaries corresponds to the number of times the classifier switches labels as the feature vector varies. We show that classifiers with fewer decision boundaries are easier to learn due to their lower Vapnik--Chervonenkis (VC) dimension.
The second aspect is classifier fairness, which refers to fair distribution of errors across classes. If in addition to classification performance, fairness is also in consideration, then the receiver can pick among the many classification rules, the one which proportionately distributes the error across classes. The third aspect is the sender's payoff under a strategic classifier. This aspect becomes relevant if a sender can opt-out and obtain a reservation payoff. Thus, a receiver may select a classifier that gives the sender a payoff higher than its reservation payoff. Although we work with scalar features, our problem suffices to demonstrate the structural aspects of strategic classification and the importance of a penalty label. Moreover, in many university admission problems like those we consider, the final admission decision has to be done based on a single scalar score, whereby the scalar problem is practically relevant.

We introduce two new algorithms called vanilla ERM and strategy-aware ERM to address the learning problem. In vanilla ERM, the classifier first uses ERM to learn the true function from the original hypothesis class. The decision problem is then applied to the learned hypothesis. In the strategy-aware ERM approach, the hypothesis class is modified to incorporate the underlying game-theoretic structure. This modification involves a two-step process. First, each hypothesis in the original class is replaced with its corresponding strategic classifier, which is obtained by applying the decision problem to each hypothesis. Second, the strategic classifier obtained in the previous step is further replaced with a new hypothesis, formed by composing the strategic classifier with the sender’s worst-case best response. The classifier then uses the Empirical Risk Minimization (ERM) algorithm to learn from this modified hypothesis class. What makes strategy-aware ERM novel is its ability to transform an otherwise unlearnable problem into a learnable one. We show that there exist hypothesis classes that are unlearnable under vanilla ERM but become learnable with strategy-aware ERM, illustrating how incorporating game-theoretic structure can fundamentally improve learnability.

While strategy-aware ERM can learn from a broader range of hypothesis classes, its advantage in the ease of learning over vanilla ERM only sometimes holds. This ease of learning is controlled by the VC dimension of the hypothesis class, which in turn determines the sample complexity in PAC learning. The comparison, therefore, becomes context-dependent if a hypothesis class is learnable via vanilla ERM. We provide a heuristic argument that strategy-aware ERM generally makes learning easier if the original hypothesis class is too complex, as it can reduce the effective VC dimension. On the other hand, vanilla ERM proves more efficient if the original hypothesis class is simple. We show this with the help of an example. A formal result, even for scalar features, remains elusive.

\subsection{Related Work}
This chapter contributes to the literature on strategic classification and learning in the presence of strategic agents~\cite{bruckner2011stackelberg,ghalme2021strategic,bechavod2022information,cohen2024bayesian,jagadeesan2021alternative,dalvi2004adversarial}, where interactions are commonly formulated as games between a strategic sender and a classifier. We model this interaction as a sender--receiver screening game and introduce a penalty label that expands the receiver's strategy space, giving rise to equilibrium strategies robust against the sender's manipulation. Preliminary versions of the results in this chapter appeared in the conference papers~\cite{singh22strategic,singh2024optimal}. This chapter is derived from our manuscript~\cite{singh2025strategic}. To the best of our knowledge, the setting of strategic classification under non-uniform preferences with penalty labels has not been previously studied, and was first explored in our work~\cite{singh2025strategic}.

Works such as \cite{hardt2016strategic,cullina2018pac,zhang2021incentive,sundaram2021pac} consider the impact of a strategic sender on the PAC learnability of the hypothesis class. In \cite{hardt2016strategic,cullina2018pac,zhang2021incentive}, the sender is assumed to have a uniform preference for positive labels, whereas in \cite{sundaram2021pac}, the sender is assumed to have a non-uniform preference for labels for each feature vector. In \cite{hardt2016strategic,cullina2018pac,sundaram2021pac}, strategic manipulation incurs a cost, whereas in \cite{zhang2021incentive}, strategic manipulation does not incur a cost. The works \cite{hardt2016strategic,cullina2018pac,zhang2021incentive,sundaram2021pac} do not explicitly utilize the underlying sender-receiver game in the learning process. The learning process in these works can be equivalently described as using the ERM algorithm to optimize over a modified hypothesis class, obtained by replacing each hypothesis with the composition of the hypothesis and the sender's best response to that hypothesis. In contrast, in this chapter, the sender has a non-uniform preference over labels for each feature vector, and the structure of this preference differs from that in \cite{sundaram2021pac}. Furthermore, in this chapter, there is a cost incurred for modifying a feature vector. We explicitly make use of the underlying sender-receiver game in the learning process by allowing the ERM algorithm to optimize over a modified hypothesis class, which is obtained by replacing each hypothesis with its corresponding Stackelberg equilibrium composed with the sender's best response.

There are other interesting formulations of the problem that we do not consider in this chapter. For instance, some research explores scenarios where the sender lacks knowledge of the classifier. Notable works in this area include \cite{ghalme2021strategic}, \cite{bechavod2022information}, and \cite{cohen2024bayesian}. Additionally, other studies introduce more structure to the sender by incorporating user types, as seen in \cite{jagadeesan2021alternative} and \cite{levanon2022generalized}.

Finally, the social welfare implications of strategic classification have been explored in works such as \cite{milli2019social}, which analyze the social costs imposed on the sender. In this chapter, we briefly address related issues such as fairness and participation constraints in Section~\ref{subsec:Other equilibria}.
